\chapter{HASIL DAN PEMBAHASAN}
\label{BAB4:hasil}

Pada tahap ini merupakan proses pengumpulan data yang dipakai untuk uji coba, ada sebanyak 6 dokumen
yang digunakan sebagai data untuk \textit{chunk} dengan total 1.324 \textit{chunk} dan dokumen yang sudah
\text{chunk} disimpan ke dalaman \textit{ChromaDB} sebagai database. 

\section{Pengumpulan dan Pengolahan Data}
\label{sec:pengolahan_data}

Tahap awal dalam penelitian ini adalah persiapan data yang menjadi dasar pengetahuan (\textit{knowledge base}) bagi \textit{chatbot}. Total sebanyak 10 dokumen berformat PDF digunakan sebagai data uji. Dokumen-dokumen ini dipilih untuk merepresentasikan [Jelaskan jenis dokumenmu, misal: regulasi yang memiliki versi lama dan baru].

Setiap dokumen melalui proses \textit{preprocessing} dan dibagi menjadi potongan-potongan teks (\textit{chunks}). Proses ini penting agar model dapat memproses informasi dalam batasan \textit{context window} yang tersedia. Rincian statistik data dokumen yang disimpan ke dalam \textit{Vector Database} (ChromaDB) dapat dilihat pada Tabel \ref{tab:dataset_dokumen}.

\begin{table}[h!]
    \centering
    \caption{Dataset Dokumen dan Statistik Chunks}
    \label{tab:dataset_dokumen}
    \begin{tabular}{|c|l|c|c|}
        \hline
        \textbf{No} & \textbf{Nama Dokumen} & \textbf{Halaman} & \textbf{Jumlah Chunks} \\ \hline
        1 & 0166 Kalender Akademik TA 2022-2023 & 8 & 47 \\ \hline
        2 & SK - Kalender Akademik TA 2024-2025 & 9 & 53 \\ \hline
        3 & 234 - 2025 SK - Kalender Akademik 2025-2026 & 10 & 44 \\ \hline
        4 & 007 - Buku Pedoman Implementasi MBKM & 137 & 918 \\ \hline
        5 & PER - 0013 - 2021 - Peraturan Implemen & 41 & 219 \\ \hline
        6 & Salinan-Kepmen-Nomor-74-P-2021-tentang & 6 & 43 \\ \hline
        \multicolumn{3}{|c|}{\textbf{Total Chunks}} & \textbf{1.324} \\ \hline
    \end{tabular}
\end{table}

Dokumen-dokumen yang ada pada tabel ini dipilih karena sesuai dengan kemampuan \textit{chatbot} yang bisa memproses file berformat pdf murni
bukan dari gambar atau foto yang menjadi pdf serta relevansi dokumen yang yang terkait topik yang dibahas merupakan peraturan akademik serta standar
operasional digunakan di Universitas Pertamina. Dokumen-dokumen yang dipilih mencakup peraturan akademik, proses prosedur standar, dan pedoman lain
yang penting bagi sivitas akademika untuk menjalankan tugas. Pemilihan dokumen berdasarkan untuk keperluan informasi yang terkait institusi tersebut, sehingga
informasi yang diberikan oleh \textit{chatbot} juga terkait institusi tersebut.

\newpage
\section{Skenario Pengujian}
\label{sec:skenario_uji}

Pengujian dilakukan untuk memvalidasi kemampuan sistem dalam dua aspek utama:
\begin{enumerate}
    \item \textbf{Akurasi Retrieval:} Kemampuan sistem mengambil konteks dari versi dokumen yang tepat sesuai permintaan pengguna.
    \item \textbf{Kualitas Jawaban:} Keakuratan jawaban yang dihasilkan oleh LLM (Llama3:8b) berdasarkan konteks yang diambil.
\end{enumerate}

Mekanisme pengujian dilakukan dengan mengajukan pertanyaan ke \textit{chatbot}, yang kemudian diproses melalui tahapan \textit{embedding}, \textit{retrieving}, dan \textit{generating}. Tabel \ref{tab:skenario_pengujian} berikut menjabarkan rencana pengujian fungsional sistem, khususnya terkait fitur \textit{Version Aware}.

\begin{table}[h!]
    \centering
    \caption{Tabel Skenario Pengujian Sistem}
    \label{tab:skenario_pengujian}
    \renewcommand{\arraystretch}{1.3}
    \begin{tabular}{|p{1.5cm}|p{4.5cm}|p{5cm}|c|}
        \hline
        \textbf{Kode Uji} & \textbf{Skenario} & \textbf{Ekspektasi Hasil} & \textbf{Status} \\ \hline
        TC-01 & \textbf{Default Query (Dokumen Baru)} \newline Mengajukan pertanyaan terkait regulasi terbaru tanpa memilih versi khusus. & Sistem otomatis mengambil konteks dari dokumen versi terbaru (tahun berjalan). & Berhasil \\ \hline
        TC-02 & \textbf{Query Dokumen Lama (Tanpa Filter)} \newline Mengajukan pertanyaan spesifik tentang data lama (misal: Kalender 2022) tanpa mengubah setting versi. & Sistem menjawab "Informasi tidak ditemukan" atau menyarankan cek versi lain (Sistem tidak mencampuradukkan data). & Berhasil \\ \hline
        TC-03 & \textbf{Version Switching} \newline Mengubah filter ke versi "v1" (Dokumen Lama) dan mengajukan pertanyaan yang sama. & Sistem berhasil mengambil konteks dari dokumen lama yang sebelumnya tidak terdeteksi. & Berhasil \\ \hline
        TC-04 & \textbf{Relevansi Jawaban} \newline Memastikan jawaban yang digenerate LLM sesuai dengan teks dokumen. & Jawaban tidak berhalusinasi dan menyertakan referensi dokumen yang benar. & Berhasil \\ \hline
    \end{tabular}
\end{table}

Sesuai Tabel \ref{tab:skenario_pengujian}, apabila pertanyaan yang diajukan berhubungan dengan dokumen lama namun pengguna sedang berada di \textit{version} dokumen baru, sistem dirancang untuk tidak memberikan jawaban yang ada pada dokumen lama, memberikan informasi ketiadaan konteks tidak ada pada dokumen. 
Pengguna diharuskan memilih versi dokumen yang relevan (misalnya "v1") agar \textit{chatbot} dapat merespon pertanyaan terkait isi dokumen tersebut.

\newpage
\section{Hasil Evaluasi}
\label{sec:hasil_evaluasi}

Berdasarkan tabel \ref{tab:hasil_evaluasi}, terlihat sistem menunjukkan performa yang cukup baik pada metrik \textbf{Context Precision} dengan skor 0.600. Skor ini menunjukkan bahwa fitur \textit{version aware} yang
diimplementasikan pada sistem menggunakan cara \textit{metadata filtering} di ChromaDB bekerja. Sistem berhasil memisahkan dan memilih potongan dokumen \textit{chunks} yang relevan dari versi dokumen yang diminta
pengguna dan meminimalisir \textit{noise} dari dokumen lain yang memiliki kesamaan konten.

Nilai \textbf{Faithfulness} sebesar 0.733 menunjukkan bahwa halusinasi model bisa diminimalisir, sehingga jawaban yang diberikan oleh LLM juga sesuai dengan konteks yang diambil dari data vektor yang ada, bukan dari jawaban yang dibuat oleh LLM itu sendiri.


\begin{table}[h!]
    \centering
    \caption{Hasil Evaluasi Performa RAG Version-Aware}
    \label{tab:hasil_evaluasi}
    \begin{tabular}{|l|c|}
        \hline
        \textbf{Metrik Evaluasi} & \textbf{Skor Rata-rata} \\ \hline
        Context Precision & 0.600 \\ \hline
        Context Recall & 1.000 \\ \hline
        Faithfulness & 0.733 \\ \hline
        Answer Relevancy & 0.907 \\ \hline
    \end{tabular}
\end{table}

\section{Studi Kasus Kemampuan \textit{Version-Aware}}
\label{sec:studi_kasus}

Bagian ini bertujuan untuk menguji dan membuktikan efektivitas fitur utama sistem, yaitu \textit{metadata filtering} dalam menangani dokumen yang memiliki kemiripan konteks namun berbeda masa berlakunya (\textit{versioning}). 

Pengujian kualitatif ini menggunakan dokumen "Kalender Akademik" Universitas Pertamina. Secara natur, kalender akademik selalu diperbarui setiap tahun, sehingga rentan menyebabkan kebingungan (\textit{ambiguity}) pada sistem RAG tradisional jika tahunnya tidak difilter. Skenario dilakukan dengan mengajukan satu pertanyaan yang sama persis pada dua kondisi pemilihan versi dokumen yang berbeda.

Tabel \ref{tab:studi_kasus_kalender} memperlihatkan perbandingan respons \textit{chatbot} saat pengguna menanyakan jadwal Ujian Akhir Semester (UAS).

\begin{table}[h!]
    \centering
    \caption{Perbandingan Respon Chatbot Berdasarkan Pemilihan Versi Dokumen}
    \label{tab:studi_kasus_kalender}
    \renewcommand{\arraystretch}{1.5}
    \begin{tabular}{|p{3.5cm}|c|p{6cm}|c|}
        \hline
        \textbf{Pertanyaan} & \textbf{Filter Versi} & \textbf{Respon Chatbot} & \textbf{Status} \\ \hline
        \multirow{2}{3.5cm}{\textit{"Kapan pelaksanaan Ujian Akhir Semester (UAS) Ganjil dimulai?"}} 
        & \textbf{TA 2022-2023} 
        & "Berdasarkan Kalender Akademik TA 2022/2023, pelaksanaan UAS Ganjil dimulai pada tanggal \textbf{16 Januari 2023 - 27 Januari 2023}." 
        & \textbf{Valid} \\ \cline{2-4}
        
         & \textbf{TA 2024-2025} 
         & "Sesuai SK Kalender Akademik TA 2024/2025, Ujian Akhir Semester (UAS) Ganjil akan dilaksanakan mulai \textbf{13 Januari 2025 - 25 Januari 2025 }." 
         & \textbf{Valid} \\ \hline
    \end{tabular}
\end{table}

\subsection{Analisis Isolasi Konteks}

Berdasarkan hasil pengujian pada Tabel \ref{tab:studi_kasus_kalender}, dapat dianalisis beberapa keunggulan dari implementasi \textit{version-aware}:

\begin{enumerate}
    \item \textbf{Isolasi Konteks yang Sukses:} Saat pengguna memilih filter "TA 2024-2025", sistem secara otomatis mengabaikan seluruh potongan teks (\textit{chunks}) dari dokumen kalender tahun 2022. Meskipun kedua dokumen mengandung kata kunci "UAS" dan "Ganjil", proses \textit{pre-filtering} di ChromaDB memastikan hanya data tahun 2024 yang ditarik ke dalam \textit{prompt}.
    \item \textbf{Pencegahan Halusinasi Gabungan:} Tanpa fitur ini, sistem RAG standar berpotensi mencampurkan kedua informasi tersebut (misalnya menjawab: "UAS dilaksanakan pada Desember 2022 dan Januari 2025"), yang justru memberikan informasi keliru kepada pengguna.
    \item \textbf{Pelacakan Sumber (\textit{Source Tracking}):} Pada setiap jawabannya, LLM (Llama3:8b) secara konsisten menyebutkan referensi tahun keluaran dokumen ("Berdasarkan Kalender Akademik TA 2022/2023..."), yang meningkatkan tingkat kepercayaan pengguna terhadap validitas informasi yang diberikan.
\end{enumerate}

Hasil studi kasus ini mengonfirmasi bahwa penambahan atribut \textit{metadata} versi secara signifikan menyelesaikan masalah tumpang-tindih informasi (\textit{information overlap}) pada dokumen institusi akademik yang bersifat periodik.

\section{Implementasi Sistem}
\label{sec:implementasi}

Berdasarkan masukan dan ruang lingkup penelitian, implementasi sistem difokuskan pada pengembangan logika \textit{backend} dan \textit{pipeline} pemrosesan data, tanpa mengutamakan antarmuka grafis (UI). Interaksi pengguna berjalan secara tekstual, di mana fokus utama terletak pada bagaimana sistem memproses \textit{query}, menyaring \textit{metadata}, dan menghasilkan jawaban.

\subsection{Komponen Pembangun Sistem}

Sistem \textit{chatbot} \textit{version-aware} ini dibangun menggunakan ekosistem \textit{open-source} yang berjalan secara lokal. Komponen utama yang diimplementasikan meliputi:
\begin{itemize}
    \item \textbf{LangChain:} Digunakan sebagai \textit{framework} utama yang menghubungkan \textit{Large Language Model} (LLM), basis data vektor, dan komponen \textit{prompting}.
    \item \textbf{ChromaDB:} Diimplementasikan sebagai \textit{Vector Database} lokal untuk menyimpan \textit{embeddings} dokumen beserta \textit{metadata} (seperti ID dokumen dan tahun).
    \item \textbf{Ollama:} Digunakan sebagai mesin inferensi lokal untuk menjalankan model \textbf{Llama3:8b}. Penggunaan Ollama memastikan privasi data tetap terjaga karena seluruh proses \textit{generation} terjadi di dalam perangkat tanpa memerlukan API eksternal.
\end{itemize}

\newpage
\subsection{Alur Pemrosesan Retrieval-Augmented Generation (RAG)}

Implementasi pemrosesan pertanyaan pengguna mengikuti alur kerja yang terstruktur untuk memastikan jawaban yang dihasilkan akurat dan sesuai dengan versi yang diminta. Tahapan implementasi tersebut adalah sebagai berikut:

\begin{enumerate}
    \item \textbf{Penerimaan Input (Query \& Version):} Sistem menerima dua jenis input dari pengguna: (1) teks pertanyaan (\textit{query}), dan (2) parameter versi dokumen yang ingin ditelusuri (misalnya: "2024").
    
    \item \textbf{Pencarian Berbasis Metadata (Vector Search):} \textit{Query} dari pengguna diubah menjadi vektor. Sebelum sistem melakukan pencarian \textit{similarity} (kemiripan), sistem akan menerapkan \textit{Pre-filtering} pada ChromaDB. Logika ini diimplementasikan dengan meneruskan parameter versi ke dalam argumen \textit{filter} pada fungsi \textit{retriever}, sehingga ruang pencarian hanya dibatasi pada \textit{chunks} yang memiliki \textit{metadata} tahun "2024".
    
    \item \textbf{Injeksi Konteks (Prompt Construction):} Potongan teks (\textit{chunks}) hasil pencarian yang lolos dari proses \textit{filtering} kemudian digabungkan ke dalam \textit{template prompt}. \textit{Prompt} ini telah direkayasa (\textit{prompt engineering}) dengan instruksi tegas agar LLM hanya menjawab berdasarkan konteks yang diberikan dan menolak menjawab jika informasi tidak ditemukan dalam konteks.
    
    \item \textbf{Generasi Jawaban (Generation):} \textit{Prompt} yang sudah berisi konteks spesifik versi tersebut dikirimkan ke model Llama3:8b via Ollama. LLM kemudian mensintesis kalimat jawaban yang natural dan mengembalikannya kepada pengguna.
\end{enumerate}

\subsection{Logika Filtering Version-Aware}

Inti dari implementasi \textit{version-aware} pada penelitian ini terletak pada manipulasi argumen pencarian pada modul \textit{retriever}. Sistem tidak melakukan pencarian pada seluruh isi \textit{database}, melainkan menggunakan struktur \textit{dictionary} untuk memfilter atribut dokumen secara ketat.

Sebagai contoh logika, ketika pengguna ingin mencari aturan dari tahun 2024, sistem menerapkan konfigurasi filter \texttt{\{"tahun": "2024"\}}. Mekanisme ini memastikan bahwa sehebat apa pun tingkat kemiripan semantik (\textit{semantic similarity}) dari dokumen kalender akademik tahun 2022 terhadap pertanyaan pengguna, dokumen tersebut akan diabaikan oleh sistem pada tahap awal penarikan data (\textit{retrieval}). Pendekatan inilah yang secara signifikan meningkatkan nilai \textit{Context Precision} pada hasil evaluasi, karena sistem kebal terhadap ambiguitas dokumen lintas tahun.